% ==============================================================================
% STAGE 4 METHODOLOGY - READY TO ADD TO Chapter 3
% ==============================================================================
% This section is fully writeable without awaiting training results.
% It describes architectural design decisions, not empirical findings.
% 
% INSERT THIS INTO: LaTeX_Thesis/chapters/03_methodology.tex
% LOCATION: After Section 1 "Overview", before existing encoder/decoder sections
% ==============================================================================

\section{Stage 4: Discriminative Key-Value Pair Extraction}
\label{sec:stage4_methodology}

Stage 4 addresses the core KVP extraction task through a discriminative learning approach combining layout-aware encoding with explicit relation modeling. This section describes the architecture, design choices, and training procedure developed for Stages 4a (entity classification) and 4b (entity + link prediction).

\subsection{Architectural Overview}
\label{sec:stage4_architecture}

The Stage 4 architecture departs from the generative approach of Stage 3 by formulating KVP extraction as a structured prediction problem with two sub-tasks:

\begin{enumerate}
    \item \textbf{Entity Recognition} (Stages 4a and 4b): Classify each token in the document as \texttt{Other}, \texttt{Key}, or \texttt{Value}.
    \item \textbf{Relation Linking} (Stage 4b only): Predict which Value tokens correspond to each Key token via scored attention.
\end{enumerate}

This decomposition offers two distinct advantages over the generative baseline:

\begin{enumerate}
    \item \textbf{Spatial grounding}: Discriminative models inherit bounding box predictions from the encoder without hallucination, whereas generative models must construct bounding boxes from text representations.
    \item \textbf{Unsupervised domain adaptation}: Transition learning can be invoked from Stage 3 (text encoding) into Stage 4 (spatial linking), leveraging the lmdx\_text representations already available.
\end{enumerate}

The encoder backbone is LayoutLMv3 (microsoft/layoutlmv3-base), chosen for its state-of-the-art multimodal fusion and proven effectiveness on form-understanding benchmarks (FUNSD, CORD, DocVQA). Stage 4a adds a simple softmax classifier over the encoder output to produce entity labels; Stage 4b adds an additional biaffine scoring module to predict relation scores.

\subsection{Why LayoutLMv3 as Encoder}
\label{sec:layoutlmv3_justification}

LayoutLMv3 was selected over alternative layout-aware encoders (LayoutLMv2, VisualBERT, Donut) for the following reasons:

\begin{enumerate}
    \item \textbf{Unified modality fusion}: Unlike LayoutLMv2 (which appends a CNN), LayoutLMv3 uses a single transformer backbone for both text and image, simplifying integration with downstream classification heads and reducing parameter count.
    
    \item \textbf{Spatial awareness without overhead}: LayoutLMv3 preserves 2D position embeddings (bounding box coordinates normalized to $[0, 1000]$ range) within its transformer layers, enabling the model to implicitly learn spatial relationships without explicit graph construction.
    
    \item \textbf{Training stability}: Pre-training on masked language modeling, masked image modeling, and word-patch alignment objectives produces robust representations, reducing fine-tuning instability compared to models pre-trained on vision-only objectives.
    
    \item \textbf{Proven downstream performance}: LayoutLMv3 achieves state-of-the-art results on FUNSD (F1 = 0.96) and CORD (F1 = 0.97) at the time of this thesis (February 2026), outperforming LayoutLMv2 and other layout-aware competitors on form understanding benchmarks.
\end{enumerate}

\subsection{Stage 4a: Entity Classification}
\label{sec:stage4a_entity_classification}

\subsubsection{Architecture}

Stage 4a is the simplest and most data-efficient variant. The LayoutLMv3 encoder output (sequence of hidden states, one per token) is fed to a linear projection layer followed by softmax:

\[
P(\text{label}_i \mid \text{token}_i) = \text{softmax}(W_{\text{class}} h_i + b_{\text{class}})
\]

where $h_i$ is the LayoutLMv3 hidden state for token $i$, and $W_{\text{class}} \in \mathbb{R}^{768 \times 3}$ is a learned projection to the three label classes (Other, Key, Value).

\subsubsection{Training Objective}

Stage 4a uses standard cross-entropy loss over all tokens:

\[
\mathcal{L}_{\text{entity}} = -\sum_{i=1}^{n} \log P(y_i^* \mid \text{token}_i)
\]

where $y_i^* \in \{0, 1, 2\}$ is the ground-truth label for token $i$, and $n$ is the sequence length (padded to 512 tokens).

\subsubsection{Label Generation from Ground Truth}

Ground-truth entity labels are generated from the KVP10k ground truth via bounding box overlap matching:

\begin{enumerate}
    \item For each key-value pair $(K, V)$ in the document's ground truth annotation, retrieve the bounding boxes $\text{bbox}_K$ and $\text{bbox}_V$.
    
    \item For each extracted word $w_i$ with bounding box $\text{bbox}_i$, compute the intersection-over-union (IoU) with both $\text{bbox}_K$ and $\text{bbox}_V$.
    
    \item If $\text{IoU}(\text{bbox}_i, \text{bbox}_K) > 0.5$, label all tokens in word $w_i$ as Key. Similarly for Value.
    
    \item All remaining tokens default to label Other.
\end{enumerate}

This bbox-based alignment avoids text-matching ambiguity and leverages the spatial grounding already present in the KVP10k annotations.

\subsection{Stage 4b: Entity + Relation Linking}
\label{sec:stage4b_entity_relation}

\subsubsection{Architecture}

Stage 4b extends Stage 4a by adding a biaffine attention module to score potential key-value relationships. After token-level entity labels are produced, the model computes pairwise relation scores:

\[
s_{i,j} = h_i^T W_{\text{biaffine}} h_j + w_{\text{bias}}
\]

where $h_i$ and $h_j$ are LayoutLMv3 hidden states for tokens $i$ and $j$, $W_{\text{biaffine}} \in \mathbb{R}^{768 \times 768}$ is a learned bilinear matrix, and $w_{\text{bias}}$ is a scalar bias term.

The relation matrix is then thresholded and masked to only allow edges from Key tokens to Value tokens based on the Stage 4b entity predictions, yielding a sparse directed graph of key-value associations.

\subsubsection{Biaffine Attention Justification}

Biaffine attention was chosen for relation scoring for the following reasons:

\begin{enumerate}
    \item \textbf{Demonstrated effectiveness}: Biaffine layers achieve state-of-the-art results in dependency parsing (Dozat \& Manning, 2017) and semantic role labeling, both requiring scoring pairwise token relations.
    
    \item \textbf{Computational efficiency}: Unlike dense cross-attention (which scales as $O(n^2)$), biaffine scoring is a simple matrix multiplication and scales effectively to documents with hundreds of tokens.
    
    \item \textbf{Interpretability}: The biaffine matrix captures symmetric and asymmetric token interaction patterns, and attention weights can be visualized to diagnose failure modes.
    
    \item \textbf{Flexibility}: The biaffine parametrization naturally supports application-specific masking (e.g., constraining links to only originate from Key tokens).
\end{enumerate}

\subsubsection{Training Objective}

Stage 4b optimizes a weighted combination of entity and relation losses:

\[
\mathcal{L}_{\text{total}} = \mathcal{L}_{\text{entity}} + \lambda \, \mathcal{L}_{\text{relation}}
\]

where $\mathcal{L}_{\text{entity}}$ is the cross-entropy loss from Stage 4a, and $\mathcal{L}_{\text{relation}}$ is binary cross-entropy over predicted relation scores:

\[
\mathcal{L}_{\text{relation}} = -\sum_{i,j} \left[ r^*_{i,j} \log \sigma(s_{i,j}) + (1 - r^*_{i,j}) \log(1 - \sigma(s_{i,j})) \right]
\]

Here $r^*_{i,j}$ is the ground-truth relation label (1 if token $i$ is a Key and token $j$ is its Value, 0 otherwise), and $\sigma$ is the sigmoid function.

\subsubsection{Hyperparameter Sweep}

To explore the interplay between entity and relation learning, we conduct an ablation study varying $\lambda \in \{0.5, 1.0, 2.0\}$:

\begin{itemize}
    \item $\lambda = 0.5$: Emphasize entity learning, relation learning is secondary
    \item $\lambda = 1.0$: Balanced weighting of entity and relation objectives
    \item $\lambda = 2.0$: Emphasize relation learning; aggressive link prediction
\end{itemize}

This sweep allows us to determine whether the task is entity-limited (relations are easy given perfect entities) or relation-limited (entities are easy but linking is hard).

\subsection{Data Preprocessing and Alignment}
\label{sec:stage4_data_prep}

The Stage 4 pipeline processes raw KVP10k data through several alignment steps:

\subsubsection{Text and Layout Extraction}

KVP10k provides lmdx\_text format: a sequence of extracted text lines, each annotated with pixel-level bounding boxes. Preprocessing extracts:

\begin{enumerate}
    \item \textbf{Word tokens}: Split each lmdx line on whitespace to obtain individual word-level text and bounding boxes.
    \item \textbf{Coordinate normalization}: Scale pixel coordinates from image-native range to $[0, 1000]$ (LayoutLMv3 standard).
    \item \textbf{Image handling}: When real document images are unavailable, create dummy 224×224 white images to satisfy LayoutLMv3Processor requirements.
\end{enumerate}

\subsubsection{LayoutLMv3Processor Configuration}

The LayoutLMv3Processor is configured with \texttt{apply\_ocr=False} to respect pre-extracted text and bounding boxes rather than performing new OCR. The processor:

\begin{enumerate}
    \item Tokenizes word sequences using the RoBERTa tokenizer (subword tokenization).
    \item Aligns subword tokens back to word-level bounding boxes via \texttt{word\_ids()} mapping.
    \item Produces input tensors: \texttt{input\_ids} (token IDs), \texttt{bbox} (normalized coordinates), \texttt{pixel\_values} (image embeddings).
\end{enumerate}

\subsubsection{Token-to-Word Label Alignment}

Ground-truth entity labels are generated at word level but must be aligned to subword tokens for loss computation:

\[
y_{\text{token}} = y_{\text{word}} \quad \text{if } \texttt{word\_ids}[\text{token}] = \text{word}
\]

For special tokens (CLS, SEP, padding) where \texttt{word\_ids} is None, labels are set to 0 (Other) to exclude from loss computation.

\subsection{Training Configuration}
\label{sec:stage4_training_config}

\subsubsection{Optimizer and Learning Rate}

Both Stage 4a and 4b use AdamW optimization with:

\begin{itemize}
    \item Initial learning rate: $\text{lr} = 2 \times 10^{-5}$
    \item Warmup steps: 1000 (5\% of training)
    \item Weight decay: $\lambda_{\text{decay}} = 0.01$
    \item Gradient clipping: $\|\nabla\| \leq 1.0$
\end{itemize}

The learning rate schedule uses linear decay after warmup, following standard transformer fine-tuning practice.

\subsubsection{Batch Size and Epoch Count}

\begin{itemize}
    \item Batch size: 16 samples per GPU (effective batch size 16 × 4 GPUs = 64 on A100 hardware)
    \item Training epochs: 5 (chosen based on typical BERT fine-tuning convergence on similar-sized tasks)
    \item Validation frequency: Every 100 mini-batches
    \item Early stopping: Patience = 3 epochs (stop if validation F1 does not improve)
\end{itemize}

\subsubsection{Hardware and Time Estimates}

Training is conducted on a single A100 GPU at (estimated):

\begin{itemize}
    \item Stage 4a: \textasciitilde 4 hours for 5 epochs over 5,389 training samples
    \item Stage 4b (single $\lambda$): \textasciitilde 5 hours per run (additional relation scoring overhead)
    \item Total for all 4b experiments ($\lambda \in \{0.5, 1.0, 2.0\}$): \textasciitilde 19 hours
\end{itemize}

\subsection{Inference and Postprocessing}
\label{sec:stage4_inference}

\subsubsection{Stage 4a Inference}

At inference, the model produces token-level probability distributions over label classes. Predictions are:

\begin{enumerate}
    \item Take argmax over predicted label probabilities to assign hard labels.
    \item Aggregate subword token labels back to word level by taking the label of the first subword token of each word.
    \item Extract contiguous spans of Key or Value tokens as final key and value phrases.
\end{enumerate}

\subsubsection{Stage 4b Inference}

Stage 4b additionally computes relation scores and constructs key-value associations:

\begin{enumerate}
    \item Predict entity labels as in Stage 4a.
    \item Compute relation scores $s_{i,j}$ for all token pairs.
    \item Extract maximum-weight bipartite matching between Key and Value spans using the relation scores as edge weights.
    \item Return (key phrase, value phrase) pairs along with confidence scores from the linker.
\end{enumerate}

\subsubsection{Bounding Box Prediction}

Both Stage 4a and 4b inherit bounding boxes directly from the preprocessed lmdx\_text, avoiding the need for separate bbox regression. This design choice:

\begin{itemize}
    \item Eliminates hallucinated bounding boxes (common in generative models)
    \item Maintains spatial grounding consistency with ground truth
    \item Simplifies the model (no bbox regression head required)
\end{itemize}

\subsection{Ablation Study Design}
\label{sec:stage4_ablation}

Stage 4 is structured around a series of controlled ablations to isolate the contributions of each component:

\begin{enumerate}
    \item \textbf{Stage 4a (entity-only) vs Stage 4b (entity + relation)}: Measures the performance gain from explicit relation learning. If Stage 4b substantially improves F1, relations are learnable; if Stage 4a is already strong, the task is entity-dominated.
    
    \item \textbf{Relation weight $\lambda$ sweep}: Stage 4b is trained with $\lambda \in \{0.5, 1.0, 2.0\}$ to determine whether the model benefits from emphasizing relation predictions ($\lambda > 1$) or entity predictions ($\lambda < 1$).
    
    \item \textbf{Stage 3 vs Stage 4 comparison}: Comparing generative Stage 3 (Mistral-7B, text-only) against discriminative Stage 4 (LayoutLMv3, layout-aware) isolates the contribution of spatial grounding to end-to-end KVP extraction accuracy.
\end{enumerate}

These ablations directly test the hypothesis that layout awareness and spatial grounding improve over text-only generative approaches on the KVP extraction task.
